\begin{theorem}[Prime-degree ramification dichotomy]\label{mod:ramification-primedegreedichotomy}
Let $K/F$ be a finite valuative extension of nonarchimedean local fields
equipped with the package from
\ref{mod:ramification-primecyclicextension}.  Thus $K/F$ is cyclic and
Galois, and its canonical degree
\[
  \ell=\operatorname{finrank}_F K
\]
is prime.  Write
\[
 \begin{aligned}
  e&=e(K/F)=\texttt{ramificationIndex}(F,K),\\
  f&=f(K/F)=\texttt{residueDegree}(F,K),\\
  p&=\operatorname{char}k_F=\texttt{residueCharacteristic}(F).
 \end{aligned}
\]
Lean uses the standard divisibility definitions
\[
 \texttt{IsTamelyRamified}(F,K)\iff p\nmid e,
 \qquad
 \texttt{IsWildlyRamified}(F,K)\iff p\mid e.
\]
In particular the tame convention includes the unramified case.

The canonical local-field identity $\ell=ef$ and primality of $\ell$
give
\[
  e=1\quad\text{or}\quad f=1.
\]
For finite extensions of local fields these are respectively the
unramified and totally ramified alternatives.  Hence if the extension is
ramified, in the precise sense $e\ne1$, then $f=1$.

Tame and wild ramification are complementary and exhaustive.  In the
wild branch, $p\mid e$ and $e\mid ef=\ell$, so $p\mid\ell$.  Both $p$ and
$\ell$ are prime, and therefore
\[
  p=\ell.
\]
Finally primality also gives the exhaustive split
\[
  \ell=2\quad\text{or}\quad \ell\text{ is odd}.
\]

Besides the individual implications, Lean exports the combined ramified
classification
\[
 f=1\quad\text{and}\quad
 \left(
   \begin{array}{l}
   K/F\text{ is tame and }(\ell=2\text{ or }\ell\text{ is odd}),\\
   \text{or}\\
   K/F\text{ is wild, }p=\ell,\text{ and }
     (\ell=2\text{ or }\ell\text{ is odd}),
   \end{array}
 \right)
\]
under the sole additional premise $e\ne1$.  The fully exhaustive dispatch
form is the disjunction of $e=1$ with that ramified classification.  No
ramification conclusion is added to the prime-cyclic structure or assumed
as a separate hypothesis.  This is exactly the case organization used in
the proof of Theorem~\ref{thm:first-main} in
\ref{sec:first-main-completion}.
\LeanFile{LanglandsFirstMainLemma/Ramification/PrimeDegreeDichotomy.lean}
\lean{LanglandsFirstMainLemma.IsTamelyRamified}
\lean{LanglandsFirstMainLemma.IsWildlyRamified}
\lean{LanglandsFirstMainLemma.unramified_or_totallyRamified}
\lean{LanglandsFirstMainLemma.totallyRamified_of_ramified}
\lean{LanglandsFirstMainLemma.isTamelyRamified_or_isWildlyRamified}
\lean{LanglandsFirstMainLemma.residueCharacteristic_eq_degree_of_isWildlyRamified}
\lean{LanglandsFirstMainLemma.degree_eq_two_or_odd}
\lean{LanglandsFirstMainLemma.ramified_primeDegree_classification}
\lean{LanglandsFirstMainLemma.primeDegree_ramification_classification}
\leanok
\SourceText{The proof of Theorem~\ref{thm:first-main}, especially
Section~\ref{sec:first-main-completion}.}
\uses{mod:ramification-primecyclicextension}
\FormalizationContract The extension invariants and residue
characteristic are the canonical quantities already supplied by the
extension and prime-cyclic APIs.  Unramified means $e=1$, totally ramified
means $f=1$, tame means $p\nmid e$, and wild means $p\mid e$.  Every
classification conclusion is derived from $\ell=ef$, primality of
$\ell$ and $p$, and elementary divisibility; none is stored in a structure
or introduced as a theorem hypothesis.
\end{theorem}
